% from ega0-19-fr.tex
\oldpage[0\textsubscript{IV-1}]{73}

le $A$-homomorphisme continu $w$ cherché en prenant le composé
$N\to N/W''\xrightarrow{h}M/V'$, le diagramme
\[
  \xymatrix{
    M\ar[r]^u\ar[d] & N\ar[d]\\
    M/V''\ar[r]_{u''} & N/W''
  }
\]
étant commutatif.
\end{proof}

Lorsque les conditions équivalentes de (19.1.5) sont remplies, nous dirons que
$u$ est \emph{formellement inversible à gauche}~; comme la condition b) de
(19.1.5) implique que $u^{-1}(W'')\subset V'$, $u$ est alors un monomorphisme
formel. La terminologie est en outre justifiée par les corollaires suivants~:

\begin{corollary}[19.1.6]
\label{0.19.1.6-fr}
S'il existe un $\widehat A$-homomorphisme continu
$s:\widehat N\to\widehat M$ tel que
$s\circ\widehat u=1_{\widehat M}$, alors $u$ est formellement inversible à
gauche.
\end{corollary}

On notera que $\widehat u$ est alors un isomorphisme topologique de
$\widehat M$ sur un \emph{facteur direct topologique} de $\widehat N$. Pour
prouver (19.1.6), il suffit de noter qu'avec les notations de (19.1.5, a)),
$v:M\to E$ se prolonge en un $\widehat A$-homomorphisme continu
$\widehat v:\widehat M\to E$ puisque $E$ est discret~; soient
$j:M\to\widehat M$ et $j':N\to\widehat N$ les homomorphismes canoniques~;
alors, en posant $w=\widehat v\circ s\circ j'$, on a
$w\circ u=\widehat v\circ s\circ\widehat u\circ j=\widehat v\circ j=v$.

\begin{corollary}[19.1.7]
\label{0.19.1.7-fr}
Supposons que les topologies de $M$ et de $N$ soient déduites de celle de $A$,
et soit $(\mathfrak J_\lambda)$ un système fondamental d'idéaux ouverts dans
$A$. Pour que $u$ soit formellement inversible à gauche il faut et il suffit
que, pour tout $\lambda$, l'homomorphisme
$u_\lambda:M/\mathfrak J_\lambda M\to N/\mathfrak J_\lambda N$, déduit de
$u$ par passage aux quotients, soit inversible à gauche (autrement dit, soit
un isomorphisme de $M/\mathfrak J_\lambda M$ sur un facteur direct de
$N/\mathfrak J_\lambda N$).
\end{corollary}

En effet, la condition est suffisante, car, pour
$V'=\mathfrak J_\lambda M$ dans la condition b) de (19.1.5), on répond à la
question en prenant $W''=\mathfrak J_\lambda N$, $V''=V'$ et $h$ tel que
$h\circ u_\lambda$ soit l'identité dans $M/\mathfrak J_\lambda M$.
Inversement, si $u$ est formellement inversible à gauche, alors, pour tout
$\lambda$, il y a, en vertu de (19.1.5, b)), un $\mu$ tel que
$\mathfrak J_\mu\subset\mathfrak J_\lambda$ et un homomorphisme
$h:N/\mathfrak J_\mu N\to M/\mathfrak J_\lambda M$ tels que l'homomorphisme
canonique $M/\mathfrak J_\mu M\to M/\mathfrak J_\lambda M$ se factorise en
\[
  M/\mathfrak J_\mu M\xrightarrow{u_\mu}N/\mathfrak J_\mu N
  \xrightarrow{h}M/\mathfrak J_\lambda M~;
\]
mais comme
$h(\mathfrak J_\lambda(N/\mathfrak J_\mu N))
\subset\mathfrak J_\lambda(M/\mathfrak J_\lambda M)=0$, $h$ se factorise
canoniquement en
\[
  N/\mathfrak J_\mu N\longrightarrow N/\mathfrak J_\lambda N
  \xrightarrow{s}M/\mathfrak J_\lambda M,
\]
et il est immédiat que $s$ est inverse à gauche de $u_\lambda$.

\begin{proposition}[19.1.8]
\label{0.19.1.8-fr}
Soit $A$ un anneau topologique admettant un système fondamental dénombrable
décroissant $(\mathfrak J_n)$ d'idéaux ouverts. Soient $M$, $N$ deux
$A$-modules topologiques dont les topologies sont déduites de celle de $A$,
$u:M\to N$ un $A$-homomorphisme. Posons
$M_n=M/\mathfrak J_nM$, $N_n=N/\mathfrak J_nN$ et soit
$u_n:M_n\to N_n$ le $(A/\mathfrak J_n)$-homomorphisme déduit de $u$ par
passage aux quotients. Supposons que, pour tout $n$,
$P_n=\operatorname{Coker}(u_n)$ soit un $(A/\mathfrak J_n)$-module projectif
et que $M$ soit séparé et complet. Alors, pour que $u$ soit formellement
inversible à gauche, il faut et il suffit que $u$ soit inversible à gauche (et
$u$ est alors un isomorphisme topologique de $M$ sur un facteur direct
topologique de $N$).
\end{proposition}

